\newpage
\begin{center}
	{ЗАКЛЮЧЕНИЕ}
\end{center}

В ходе выполнения выпускной квалификационной работы была разработана информационная система на платформе «1С:Предприятие 8.3» для комплексного учета деятельности спортивной организации. Разработанная конфигурация позволяет автоматизировать ключевые процессы спортивной организации и обеспечивает удобный доступ к информации для различных категорий пользователей.

На этапе анализа предметной области была охарактеризована деятельность спортивной организации на примере спортивной школы, выявлены основные бизнес-процессы, требующие автоматизации, и проведен обзор существующих программных решений. Результатом анализа стало обоснование необходимости разработки собственной конфигурации, полностью адаптированной под нужды спортивного учреждения.

В соответствии с техническим заданием в системе реализованы все заявленные функции:

\begin{enumerate}
	\item Управление данными о спортсменах и тренерах: справочники позволяют вести анкетные данные, фиксировать виды спорта и присвоенные разряды для спортсменов, а также специализацию и должность для тренеров.
	\item Составление и ведение расписания занятий: документ «Расписание занятий» обеспечивает планирование тренировок на месяц с детализацией по дням недели, тренерам, группам и местам проведения.
	\item Управление соревнованиями: документ «Соревнования» позволяет регистрировать спортивные мероприятия с указанием трех табличных частей: дисциплины, участники и судьи.
	\item Учет результатов соревнований: документ «Результаты соревнований», создаваемый на основании соревнований, обеспечивает ввод и хранение результатов с возможностью загрузки из файлов Excel.
	\item Оформление приказов на присвоение разрядов: «Приказ о присвоении разряда» автоматизирует процесс присвоения спортивных разрядов, имеет печатную форму.
	\item Складской учет спортивного инвентаря: четыре складских документа (поступление, выдача, возврат, списание) формируют движения по регистру накопления, позволяя получать актуальные остатки в любой момент.
	\item Формирование отчетности: реализованы отчеты «Расписание занятий общее», «Расписание занятий группы», «Расписание занятий тренера», «Список спортсменов» и «Результаты соревнований».
	\item Разграничение доступа: система поддерживает две роли: администратор (полный доступ) и тренер (ограниченный доступ без права работы со складскими документами, приказами и настройками).
	\item Отправка уведомлений: реализована рассылка уведомлений участникам соревнований по электронной почте.
\end{enumerate}

Созданная конфигурация включает 21 справочник, 10 документов, 4 перечисления, 1 регистр накопления и 4 регистра сведений. Интерфейс системы разделен на четыре подсистемы в соответствии с требованиями технического задания: «Учет спортсменов и тренеров», «Учет соревнований», «Учет спортивного инвентаря» и «Администрирование». Все формы выполнены в едином стиле управляемого приложения.

Практическая значимость работы заключается в создании готового к внедрению программного продукта, позволяющего сократить временные затраты на ведение документации, повысить качество планирования тренировочного процесса и обеспечить руководство спортивной организации актуальной информацией для принятия управленческих решений.

Дальнейшее развитие системы может быть направлено на расширение функционала: добавление модуля учета посещаемости, реализацию личного кабинета спортсмена с доступом через веб-интерфейс, интеграцию с мобильными приложениями, а также внедрение автоматизированного тестирования на базе фреймворка Vanessa Automation.

Таким образом, поставленные цели и задачи выпускной квалификационной работы выполнены в полном объеме. Разработанная информационная система готова к внедрению в деятельность спортивной организации.